\lecture{4}{Fri 02 Sep 2022 21:27}{Continuity and Differentiability}

\subsection{Continuity and Differentiability}

\begin{definition}[Continuity]
	$f$ is continuous at a point $a$ in the domain if \[
	\lim_{z \to a} f(z) = f(a)
	.\] 
\end{definition}

\begin{example}[Continuous functions]
	\quad
	\begin{itemize}
		\item $f(z) = c$ 
		\item $f(z) = z$
		\item  $fg, f+g$ are continuous when $f$ and $g$ are continuous
		\item $\frac{f}{g}$ continous when $f, g$ continous and $g(z) \neq  0$

	\end{itemize}
	
\end{example}

\begin{proposition}
	 
	$\lim_{z \to a} f(z) = f(a)  \iff u, v$ are continuous at $(a_1, a_2)$ where $f = u + iv $ and $a = a_1 + i a_2$	
\end{proposition}

Let $a \in \partial \mathcal{D}$ ,where the $\mathcal{D}$ is domain of $f$. If $\lim_{z \to a} f(z)= b$, define $f(a) = b$. The extended function $\overline{f} : \overline{\mathcal{D}} \to \mathbb{C}$ is continous by definition.

\begin{example}[Removable Singularity]
	$f(x) = \frac{z^2 - 1}{z_-1}$ is defined for $z \neq  1$. The function $z+1$ is defined everywhere, and continous. Such a sigularity is removable.
\end{example}
\begin{example}[Nonremovable Singularity]
	$u(x,y) = \frac{x+y}{x-y}$, $x \neq y$. The limit $\lim_{x,y \to 0} u(x,y)$ does not exist. Verified by choosing different paramaterizations and yield different limits.

	$\lim_{z \to 0} \frac{\overline{z} }{z}$ does not exist.
\end{example}

\begin{remark}
	It only makes sense to put approximate value into \textit{continuous} functions, since a little change in input should result in a little change in result.
	
\end{remark}

\begin{definition}
	Let $f$ be a complex function defined in a region (open, connected set) $\mathcal{D}$. We say that $f$ is $\mathbb{C}$-\textit{differentiable} at $a \in \mathcal{D}$ if \[
	\lim_{h \to 0} \frac{f(a+h)-f(a)}{a}
	\] 
	exists. We define $f'(a)$ to be the limit.
\end{definition}

\begin{remark}
	Differentiability in real and imaginary parts do not imply $\mathbb{C}$-differentiability in general.
\end{remark}

\begin{theorem}
	If $f, g$ are differentiable, then $(f+g)' = f' + g'$. Also, $(fg)' = f'g + fg'$.
\end{theorem}

\begin{definition}
	$d = \operatorname{deg} P$ is the degree of polynomial $P$ if $d$ is the largerst natural number such that $a_d \neq  0$.
\end{definition}

Formal consequences:
\begin{enumerate}
	\item Constant functions are differentiable.
	\item Identity function is differentiable, $f'(z) = 1$.
	\item Polynomials are differentiable. (using the theorem) $(z^n)' = n z^{n^{-1} }$
\end{enumerate}

\begin{theorem}[factoring, FTA]
	$P(z_0) = 0 \implies P(z) = Q(z) (z-z_0)$. A polynomial of degree $d$ has at most $d$ distinct zeros.
\end{theorem}


\begin{definition}
	Ratios of polynomials are called \textit{rational functions}.
	
\end{definition}

Singularity of a rational function is called a  \textit{pole}. The limit of rational function at a pole is infinity. Rational functions are differentiable everywhere except finitely many point, which are its poles.

Consequence of differentiability

\begin{remark}
	Differentiability does not imply smoothness. For example $f(z) = \overline{z} $ is not differentiable at $0$. Differentiability in complex analysis implies something very different from real analysis.
\end{remark}
